\documentclass[11pt]{article}
\usepackage[margin=1in]{geometry}
\usepackage{amsmath,amssymb,booktabs}
\usepackage[hidelinks]{hyperref}
\setlength{\parindent}{0pt}
\setlength{\parskip}{.5em}
\newcommand{\ind}{\mathbf{1}}
\newcommand{\code}[1]{\texttt{\detokenize{#1}}}
\title{Optional Six-Manager Evacuation Experiment\\Mathematical Construction, Placement Search and Initial Results}
\author{}
\date{}
\begin{document}
\maketitle

\begin{abstract}
This document specifies the implemented optional crowd-manager model and
summarizes the saved experiments with 300 audience members. Six stationary
managers communicate shortest-distance routes with proximity-dependent
authority, while guests retain their existing behavioral preferences. Physical
closures and knowledge of closures are modeled separately. A discrete genetic
search evaluates placements through repeated microscopic simulations. In the
ordinary-policy reproduction, automatic placement reduced total clearance from
60.443 to 58.742 seconds. With counter-keyed paired randomness, the corresponding
initial-layout comparison was 58.638 versus 58.188 seconds. Small placement
searches improved training performance, but a placement selected across blockage
scenarios was slower on a held-out audience. These are preliminary computational
results, without statistical or global-optimality guarantees.
The numerical results were obtained before the assigned-manager-priority
revision below; they describe the earlier strongest-signal selection rule.
\end{abstract}

\section{Experimental scope and baseline}
The existing model remains the default. Setting
\code{crowdManagers.enabled=true} enables six stationary managers; omitting
this object or setting its flag to false adds no staff, route graph, guidance
utility, or agent guidance records. Closure-only comparisons may explicitly
retain scheduled closures while disabling staff. The original four-dimensional
behavioral weights, their covariance, and the three-second default decision
interval remain unchanged.
The separate student-audience option overrides this experiment: with
\code{studentMode=true}, no crowd managers are active even if their stored
flag is enabled. Students directly minimize graph distance. Physical closure
schedules can remain in effect without staff or notifications. See
\code{STUDENT_MODE.tex}; the historical results here concern the ordinary
behavioral audience, not student mode.

There are $n$ audience members and exactly six additional staff discs.
The verification runs use $n=300$. Staff do not contribute to audience
population, completion counts, peer shares, or stair queues. Their stationary
bodies do participate in physical collision exclusion and interpersonal forces.
Geometry, velocities and times retain the units in
\code{BUILDING_SPECIFICATION.tex}. Geometry is schematic, not surveyed.

\subsection{Existing behavioral utility}
An action $a$ identifies a next local target. Let
$\mathcal A_i^{\mathrm{phys}}(t)$ denote physically available actions and
$\mathcal A_i^{\mathrm{obs}}(t)$ those currently visible, remembered, instructed,
or the sole permitted aisle continuation. The latter set is used for choice.
Manager directions can therefore change both information and utility.

The behavioral component is
\[
 B_i(a,t)=-w_{Ti}T_i(a,t)+w_{Pi}P_i(a,t)
                   -w_{Ni}N_i(a,t)-w_{Si}C_i(a,t).
\]
Here $T_i$ is perceived normalized travel time, $P_i$ the fraction of observed
nearby people with aligned intended direction, $N_i$ the maximum of the
counterflow and queue-cutting cues, and $C_i$ the indicator for abandoning an
existing target in the current stage. Familiarity is absent.

In the present configuration,
\[
 T_i=\frac{d_i^{\mathrm{loc}}/\max(0.35,v_i^0)
                  +\widehat W_i^{\mathrm{queue}}+\widehat W_i^{\mathrm{row}}}
                 {t_0},\qquad t_0=1\ \mathrm{s},
\]
with distances in metres and speeds in metres per second.
$d_i^{\mathrm{loc}}$ is the existing policy's local-target distance, which is
distinct from the manager graph's downstream distance and the manager--guest
communication distance. Queue delays use observed older waiters ahead; the
row-delay cue reflects experienced blockage by opposing row movement.

The four weights are sampled once per guest:
\[
 \mathbf z_i\sim\mathcal N_4(-\tfrac12\mathbf 1,\Sigma),\qquad
 \mathbf w_i=\exp(\mathbf z_i),\qquad
 \Sigma=(1-\rho)I_4+\rho\mathbf 1\mathbf 1^{\mathsf T},
 \quad\rho=-0.15.
\]
Exponentiation is componentwise. Each marginal latent variance is one, so
$\mathbb E[w_{ji}]=1$; this does not impose a fixed sum for each person.

With the default \code{behavioralOnly=false}, the complete pre-guidance utility
also contains the existing distance and stochastic terms:
\begin{align}
 U_i^{\mathrm{existing}}(a,t)
 &=D_i(a,t)+B_i(a,t)+\Pi_i(a,t)+\epsilon_i(a,t),\\
 D_i(a,t)&=\left(\frac{d_{i,\min}^{\mathrm{loc}}+d_0}
                            {d_i^{\mathrm{loc}}+d_0}\right)^{\beta_i},
             \qquad d_0=1\ \mathrm{m},\\
 \Pi_i(a,t)&=p_i\left(1-e^{-t/(60\,\mathrm{s})}\right)\delta_i(a,t).
\end{align}
Here $\beta_i\in[0,1]$ is the sampled distance sensitivity and $p_i\in[0,1]$
the sampled panic sensitivity. The panic and residual shocks have marginal
distribution $\mathcal N(0,0.2^2)$ and are held per option in 15-second buckets.
Given an instruction, the actual decision is
\[
 a_i^*(t)\in\arg\max_{a\in\mathcal A_i^{\mathrm{obs}}(t)}
                 \{U_i^{\mathrm{existing}}(a,t)+G_i(a,t)\}.
\]
Thus shortest-distance guidance is a recommendation within a heterogeneous
utility model; it is not a deterministic assignment of all guests to routes.

\subsection{Coupling to physical movement}
The selected action supplies the existing legal local waypoint. In the
manuscript force mode the acceleration contains
\[
 \dot{\mathbf v}_i=
 \frac{v_i^0\widehat{\mathbf e}_i-\mathbf v_i}{\tau}
 +\mathbf f_i^{\mathrm{obs}}
 +\sum_{j\in\mathcal N_i^{\mathrm{aud}}}\mathbf f_{ij}^{\mathrm{int}}
 +\sum_{m\in\mathcal N_i^{\mathrm{staff}}}\mathbf f_{im}^{\mathrm{int}},
 \qquad \tau=0.65\ \mathrm{s}.
\]
Staff have $\mathbf v_m=\mathbf 0$ and use the same interpersonal interaction
law within the 2.6-metre cutoff. They are excluded from the behavioral peer
denominator. Hard-body tests additionally require
\[
 \|\mathbf x_i-\mathbf y_m\|^2
       \geq 1.005\,(r_i+r_M)^2
\]
on the same plane. Euler integration, speed limits, region constraints and the
optional local collision-avoidance layer are retained. Guidance supplies no
extra physical force pointing toward the manager's body.

\section{Directed distance graph}
Let $G=(V,A)$ be the graph of declared movement states. A vertex contains a
floor, region, position, and where necessary a seating zone. Vertices represent
portal departures, named destination entrances, radial aisle ends, and their
circulation-region counterparts. Coordinate overlap never creates an edge
between distinct floors or regions. Elevator departures are excluded.

An edge $a=(u,v)$ has nonnegative length $L_a$ in metres. Its cost contains no
crowd delay, speed, utility, or capacity coefficient. Ordinary portal transfers
use the distance between declared departure and landing. Stair and access
transfers use their declared physical traversal lengths, including all flights.
Radial aisle edges point outward, reflecting the implemented action set.
Circulation edges connect pairs in the same declared region.

\subsection{Within-region distance}
In a convex rectangular region the cost is Euclidean distance. In an annular
sector with inner centre-boundary radius $R$, put
\[
 r_j=\|\mathbf x_j-\mathbf c\|,\quad
 \Delta=|\theta_1-\theta_2|,\quad
 \alpha_j=\arccos(R/r_j).
\]
For the current sectors, whose angular width is at most $\pi$,
\[
 d(\mathbf x_1,\mathbf x_2)=
 \begin{cases}
 \|\mathbf x_1-\mathbf x_2\|,&\Delta\leq\alpha_1+\alpha_2,\\
 \sqrt{r_1^2-R^2}+\sqrt{r_2^2-R^2}
       +R(\Delta-\alpha_1-\alpha_2),&\text{otherwise}.
 \end{cases}
\]
The second case joins two tangents by an inner-boundary arc. Both paths stay
inside the outer disk. The mechanical annular bounds are already centre
bounds; no additional body erosion is applied to this expression. Points
within $10^{-10}\max(1,R)$ of the inner boundary are snapped to $R$ to avoid
square-root/inverse-cosine amplification of coordinate roundoff.

This is a reference distance in the reduced schematic graph. Actual movement
uses the existing legal waypoints and social-force integrator; its realized
length and time can exceed the graph reference. The graph is not a navigation
mesh resolving every architectural detail.

\subsection{Rows and seat-to-exit routes}
For each legal row-side action, the attachment length is the remaining curved
row distance to its aisle edge, one aisle half-width to the aisle centre, and
the outward radial distance to the aisle's declared end. This explicit
row-to-aisle abstraction avoids a straight line through seating blocks.
The small mechanical row-to-aisle repositioning is therefore represented by
a finite reference length even though it is not continuously integrated.

Let $B$ be a set of known closed portal IDs. Transfer edges for $B$ and final
threshold seeds in $B$ are removed. The reverse multi-source Dijkstra field is
\[
 D_B(v)=\min_{\pi:v\leadsto V_{\mathrm{exit}}(B)}\sum_{a\in\pi}L_a,
\]
where $V_{\mathrm{exit}}(B)$ contains only exterior-threshold vertices whose
portal IDs are not in $B$, and paths use the remaining transfer edges.
For a proposed local action $a$, compute
\[
 Q_B(i,a)=L_{i,a}+D_B(v_a).
\]
For a portal action, $L_{i,a}$ includes the chosen portal's transfer and
$v_a$ is its landing: visiting a portal marker and walking past it cannot
receive credit for using that portal. An exterior action ends at its threshold.
The recommendation is a minimizer of $Q_B$ over physical local actions,
including actions not previously known to the guest. Exact ties use stable
action/exit IDs; Dijkstra uses deterministic traversal with a $10^{-10}$ metre
strict-improvement tolerance to avoid cycles at zero-length links.

Every seat receives an initial all-open reference route. Later instructions
are recomputed from current position and current knowledge. Unreachable states
produce a null route, never an invented escape or non-finite JSON value.
At most 64 knowledge-specific fields are retained in an LRU cache.

\section{Placement domain and responsibility}
Generate a finite candidate set $C$ in aisles, circulation bands, and foyer.
Each site is checked against region legality, all generated seat centres, and
portal approach clearances. The current candidates reserve at least
$0.85+r_M$ metres from portal centres and $r_M+0.2025+0.08$ metres from seat
centres. These are experimental placement restrictions, not measured rules.
The default world yields 110 candidates and an 85-vertex route graph.

Choose $M\subseteq C$ with $|M|=6$. Staff pairs on the same floor must be at
least $2r_M+0.08$ metres apart. By default, both occupied floors must have at
least one manager; the remaining allocation is optimized. This constraint can
be disabled explicitly. Candidate generation is placement discretization,
not a discretization of continuous audience motion.

Every seat is assigned a responsible manager using strongest initial signal,
then same-floor proximity and a stable identifier to break ties. An uncovered
seat still has an assignment, which does not imply it can hear the manager.
The assignment remains fixed throughout the run and has priority when several
managers can communicate with that person. Other managers can provide fallback
instructions when the assigned manager cannot be heard.

\section{Signal, utility and remembered instructions}
For audience member $i$ and manager $m$, define
\[
 q_{im}(t)=\ind\{\ell_i=\ell_m\}
 \ind\{d_{im}\leq R_c\}\ind\{\text{line of sight}\}
 \exp(-d_{im}/\lambda).
\]
Here $d_{im}$ is Euclidean distance on the same plan. The existing schematic
walls occlude communication. Communication is omnidirectional; a guest's
visual field need not point toward the manager. No cross-floor audio, acoustic
propagation, speech-service capacity, or communication delay is modeled.

Let $m_i^0$ be the manager initially assigned to guest $i$, and let
$\mathcal H_i(t)=\{m:q_{im}(t)>0\}$ contain the managers whose instructions
can currently be received. The selected manager is
\[
 m_i^*(t)=
 \begin{cases}
 m_i^0,&m_i^0\in\mathcal H_i(t),\\
 \displaystyle\arg\max_{m\in\mathcal H_i(t)}q_{im}(t),
       &m_i^0\notin\mathcal H_i(t),\ \mathcal H_i(t)\ne\varnothing.
 \end{cases}
\]
Fallback signal ties use a stable manager identifier. When $\mathcal H_i(t)$
is empty, no new instruction is received; an existing valid remembered
instruction may remain active. Receiving fallback guidance does not reassign
$m_i^0$. If the original manager becomes receivable again, it regains priority.
This preference does not bypass floor, range or wall-occlusion restrictions,
and does not guarantee obedience to the selected instruction.

Only the selected manager supplies a new instruction; signals do not add.
The utility bonus uses that manager's actual signal, even if another manager
has a stronger signal. In the original experiments reported later, selection
instead always used the strongest receivable signal; their results have not
been recomputed under assigned-manager priority.
Managers use the building graph and their own closure knowledge. On receiving
an instruction, the guest learns the recommended target and portal sequence.
Let $t_r$ be receipt time, $q_r$ its received signal, and $\tau_M$ the memory
horizon. While its route is valid and $0\leq t-t_r\leq\tau_M$, the bonus is
\[
 G_i(a,t)=\kappa\gamma_iq_r e^{-(t-t_r)/\tau_M}
             \ind\{a=a_i^{\mathrm{instruction}}(t)\}.
\]
The implemented utility is
\[
 U_i^{\mathrm{guided}}(a,t)=U_i^{\mathrm{existing}}(a,t)+G_i(a,t).
\]
The existing utility includes its optional distance, panic and residual terms;
with \code{behavioralOnly=true}, it consists of the original four-factor dot
product plus guidance. Guidance does not restore familiarity. Initially
$\gamma_i=1$ for everyone, without an additional random draw.

At a region transition outside communication range, a remembered instruction
advances only to an eligible portal already in the received sequence, or the
single physical aisle continuation. It does not run a new global search.
New instructions refresh memory. Expiry or learning that the stored route
contains a blocked portal invalidates it. Expiry removes the authority bonus;
previously learned target locations remain known. The bonus is finite, so
unbounded behavioral weights can still lead to noncompliance.

\begin{center}
\begin{tabular}{lll}
\toprule
Parameter & Default & Meaning\\
\midrule
$\kappa$ & 12 & Maximum utility bonus\\
$\lambda$ & 8 m & Spatial decay length\\
$R_c$ & 12 m & Communication range\\
$\tau_M$ & 60 s & Memory horizon and decay scale\\
$r_M$ & 0.2025 m & Stationary staff body radius\\
$\gamma_i$ & 1 & Guest responsiveness\\
\bottomrule
\end{tabular}
\end{center}
These are transparent experimental defaults, not fitted behavioral estimates.
Hold them fixed during placement optimization; use separate sensitivity studies.

\section{Physical closures and knowledge}
Each permanent closure specifies portal ID, onset $t_b$, and optional staff
notification delay $\delta$. The portal becomes unavailable at the first tick
with $t\geq t_b$, before decisions and motion; it cannot be crossed even if
everyone still believes it open. Tick-level onset is conservative within at
most one integration step. A null delay means no automatic broadcast.

Managers learn the closure from scheduled notification at $t_b+\delta$,
their own local observation within the existing discovery radius, or a report
from a guest receiving their instruction. A guest receives the heard manager's
knowledge and also retains its original local discovery mechanism. A manager
can therefore recommend an actually closed route while still uninformed.
Writing $B^{\mathrm{true}}(t)$, $B_i(t)$ and $B_m(t)$ for physical, guest and
manager closure sets makes this distinction explicit. During receipt from
manager $m$, both parties exchange their known closures, producing
$B_i\leftarrow B_i\cup B_m$ and $B_m\leftarrow B_i\cup B_m$; the route query
uses this combined knowledge. There is no query of $B^{\mathrm{true}}$ inside
the shortest-path algorithm. Valid information can arrive through observation,
notification or reporting, so knowledge need not remain absent until the
scheduled notification time.
Manager-to-manager broadcasts are not simulated. Reports are processed in the
engine's seeded sequential order; scheduled news is applied at tick boundaries.
Reopening and partially reduced capacity are not implemented by this schedule.
Closing one of five foyer doors differs from closing all five; each has its own
portal ID. To close a bank, supply an event for every affected portal.

\section{Simulation-based placement optimization}
First initialize with greedy marginal coverage
\[
 H(M)=\sum_{i\in\text{all seats}}\max_{m\in M}q_{im}(0).
\]
This is only an initialization proxy. Add random feasible layouts to obtain a
genetic population. A chromosome is a sorted tuple of six distinct candidate
IDs; permutations are one layout. Tournament selection, crossover from the
parent site union, one or two replacement mutations, feasibility rejection,
and one elite define each new generation. A separate seeded search RNG never
changes the audience RNG. Cached layouts are simulated once per training set.

Let $s$ index closure scenarios with positive weight $p_s$ and $r$ index equal
weight training audience seeds. Normalize $w_{sr}=p_s/(R\sum_s p_s)$.
With horizon $H$, remaining audience $N_{sr}(M)$, and true clearance time
$T_{sr}(M)$ only when all exit, minimize the lexicographic vector
\[
 J(M)=\left(
 \sum_{sr}w_{sr}\ind\{N_{sr}>0\},\quad
 \sum_{sr}w_{sr}\frac{N_{sr}}{n},\quad
 \sum_{sr}w_{sr}\widetilde T_{sr}\right),
\]
where $\widetilde T_{sr}=T_{sr}$ on completion and $H$ otherwise.
Thus incomplete runs are retained and rank behind complete runs. When all
complete, the third component is mean total clearance time. A capped time is
not reported as actual clearance. Scenario weights are user assumptions; equal
example weights are not empirical blockage probabilities.

The optimizer evaluates $P$ initial layouts and at most $(P-1)$ new layouts per
subsequent generation; with $K$ generations, $S$ scenarios and $R$ training
seeds the bound is $[P+K(P-1)]SR$ training simulations before cache reuse.
Held-out evaluation adds two runs per validation seed/scenario: no managers
and the selected six-manager layout. Validation seeds must be disjoint from
training seeds and never select a layout. This heuristic returns the best found
placement on $C$, with no global or continuous-domain optimality guarantee.

\section{Paired randomness and diagnostics}
Audience initialization retains the original seeded RNG. The optimizer enables
an optional counter-based shock stream keyed by
\code{(seed,agentId,noiseBucket,kind,actionId)} using SHA-256 and Box--Muller.
Hence a changed action set cannot shift another person's residual or the
mechanical shuffle RNG. A fingerprint verifies matching seats, traits, weights,
panic and distance sensitivity across baseline/guided validation pairs. Ordinary
runs retain the original stream unless \code{pairedRandomness=true} is supplied.

Reports contain training objectives, search history, settings, layouts, scenario
definitions, seed sets, matched fingerprints, held-out paired differences,
$T_{95}$, completion status, exit usage, adherence and observed queue peaks.
Adherence includes single-action continuation stages; it is an implementation
diagnostic, not an independently estimated probability of obeying staff.
For each decision, the implementation also records the winner after subtracting
only the guidance bonus from the same offered actions. A difference increments
\code{changedChoiceDecisions}. This local comparison retains the crowd state,
information and shocks; it does not estimate a full no-manager causal outcome.
The queue proxy counts slow agents with a committed queue target within the
existing observation distance, sampled at tick boundaries. It is not pressure,
hazard exposure, or a calibrated safety-density measure. Null values identify
unattained quantiles and incomplete clearance. JSON uses no infinity/NaN.

\section{Configuration used in the saved experiments}
All results below use the two-plan schematic auditorium, 300 audience members
and, when enabled, six additional stationary managers. The audience contains
119 orchestra, 110 loge and 71 balcony guests. Completion requires crossing a
first-floor exterior threshold; an interior transfer does not count as exit.

\begin{center}
\begin{tabular}{ll}
\toprule
Setting & Experimental value\\
\midrule
Integration step & $\Delta t=0.1$ s\\
Decision interval & 3 s, with immediate stage/closure reconsideration\\
Verification/search horizon & 120 s; all reported runs completed\\
Latent weight standard deviation/correlation & $\sigma=1$, $\rho=-0.15$\\
Time normalization & $t_0=1$ s\\
Shock standard deviation/bucket & 0.2 utility units / 15 s\\
Preferred speed & Clipped normal: $\mu=1.29$, $s=0.19$ m/s\\
Preferred speed bounds & 0.5--2.3 m/s\\
Reaction delay & Clipped lognormal: nominal mean 2 s, SD 1.2 s\\
Reaction delay bounds & 0--40 s\\
Guest and staff radii & 0.2025 m\\
Physical constraints & Regions, hard collisions and door capacity enabled\\
Force/behavior mode & Manuscript forces; local behavioral actions\\
Local avoidance & Enabled\\
Manager parameters & $\kappa=12$, $\lambda=8$ m, $R_c=12$ m, $\tau_M=60$ s\\
Responsiveness/floor coverage & $\gamma_i=1$; at least one manager per floor\\
\bottomrule
\end{tabular}
\end{center}
Speed and reaction parameters describe the distributions before clipping,
rather than the exact moments after clipping. The normal GUI world has a longer
default time limit, but the reproduction used 120 seconds. All compared runs
finished before this bound, so it did not censor the results.

The automatic coverage initializer selected the following six sites:
\begin{center}
\begin{tabular}{lll}
\toprule
Section & Site 1 & Site 2\\
\midrule
Orchestra & \code{aisle:orchestra:2:0} & \code{aisle:orchestra:4:1}\\
Loge & \code{aisle:loge:1:0} & \code{aisle:loge:4:0}\\
Balcony & \code{aisle:balcony:2:0} & \code{aisle:balcony:4:0}\\
\bottomrule
\end{tabular}
\end{center}
These positions depend on all generated seats, not random occupancy. They are
coverage-initialized positions, not the result of the genetic search.

\section{Initial numerical results}
All tables in this section are historical results from the original
strongest-signal selection policy. The assigned-manager-priority revision
changes manager choice when communication regions overlap; the previous
timings and selected placements are not estimates of its current performance.
\subsection{Metrics and experimental protocols}
For completed exit times $t_{(1)}\leq\cdots\leq t_{(n)}$, define
\[
 T_{100}=t_{(n)},\qquad T_{95}=t_{(\lceil0.95n\rceil)},\qquad
 \overline T=\frac1n\sum_{i=1}^{n}t_i.
\]
At $n=300$, $T_{95}$ is the 285th exit time. Let
$\Delta T=T_{100}^{\mathrm{guided}}-T_{100}^{\mathrm{baseline}}$; negative
values indicate faster clearance. Percentage change is
$100\Delta T/T_{100}^{\mathrm{baseline}}$.

Ordinary and counter-keyed paired randomness are different experimental
protocols and should not be pooled as replications of one treatment. With the
ordinary stream, changed route/option sets can shift subsequent random draws
despite identical initialization seeds. The paired protocol controls that
source of variation while permitting crowd interactions and queues to evolve
differently because of managers.

\subsection{Ordinary GUI configuration: seed 2026}
The headless engine reproduction of the ordinary GUI configuration gave:
\begin{center}
\begin{tabular}{lrrr}
\toprule
Mode & $T_{100}$ (s) & $T_{95}$ (s) & $\overline T$ (s)\\
\midrule
Managers off & 60.443 & 48.149 & 29.409\\
Six managers on & 58.742 & 48.164 & 29.218\\
\bottomrule
\end{tabular}
\end{center}
Total clearance decreased by 1.701 seconds (2.815\%), but $T_{95}$ was almost
unchanged and slightly larger with managers. This realization therefore does
not demonstrate a uniform improvement across all metrics.

The guided run recorded 2,763 decisions with a positive instruction bonus.
Subtracting only that bonus while retaining the same offered actions changed
the winning action in 137 decisions. A total of 299 guests had different exit
times. Both runs ended with guest 267, seat \code{B-07-2-5}, crossing
\code{M4}; that guest had no recorded local choice reversal attributable to
the bonus. Exit-time differences can include interaction effects and ordinary
random-stream changes; they are not all direct manager-induced choice reversals.

The GUI tick clock rounds to tenths of a second and advances to the end of a
tick. Its precise clearance display uses the interpolated final crossing time.
These are different measurements.

\subsection{Paired initial-placement comparisons: seed 2026}
With counter-keyed paired randomness and the same automatic placement:
\begin{center}
\begin{tabular}{lrrr}
\toprule
Mode & $T_{100}$ (s) & $T_{95}$ (s) & $\overline T$ (s)\\
\midrule
No managers; all exits open & 58.638 & 47.666 & 28.930\\
Six managers; all exits open & 58.188 & 46.217 & 28.692\\
Six managers; delayed blockage news & 79.699 & 66.238 & 34.327\\
\bottomrule
\end{tabular}
\end{center}
The open-exit manager benefit was 0.450 seconds, approximately 0.767\% of
baseline clearance. In the blockage case, \code{SIDE_E1} closed at 10 s and
automatic staff news was scheduled at 25 s. The guided blocked run was 21.511
seconds slower than the guided open run. This difference combines capacity
loss, rerouting and information effects; it does not isolate news delay alone.

All three runs evacuated 300 guests. The paired open baseline sent 133 guests
through the five-door \code{MAIN} bank, whereas the guided run sent 48 there.
A small change in total clearance can therefore coexist with substantial
redistribution of exit use.

\subsection{Open-exit placement search}
The first search used three initial layouts, one evolutionary generation,
training seed 2026 and held-out seed 7. Five distinct layouts were evaluated
in training. The selected positions were
\begin{center}
\begin{tabular}{ll}
\code{aisle:loge:0:0} & \code{aisle:loge:1:1}\\
\code{aisle:orchestra:0:0} & \code{balcony-rim:0:0}\\
\code{corridor:1:1} & \code{foyer:0:0}
\end{tabular}
\end{center}
Its training clearance was 54.749 seconds, versus 58.188 seconds for automatic
coverage on the same training audience. The best member of the initial search
population cleared in 54.921 seconds; the evolutionary generation improved
that value to 54.749 seconds. These are training comparisons.

On held-out seed 7, baseline clearance was 56.900 seconds and selected-layout
clearance was 56.647 seconds: a 0.253-second reduction (0.445\%). The
corresponding $T_{95}$ values were 46.957 and 46.449 seconds. The held-out
benefit was considerably smaller than the apparent training gain.

\subsection{Three-scenario layout screening}
A separate trial compared two initial layouts with zero evolutionary
generations. This was a small layout screening, not a deep optimization run.
Training seed 2026 and three equal-weight scenarios were used: all exits open,
\code{SIDE_E1} closing at 10 s with immediate staff news, and the same closure
with staff news delayed by 15 s. The selected layout was
\begin{center}
\begin{tabular}{ll}
\code{aisle:loge:0:0} & \code{aisle:loge:1:1}\\
\code{aisle:orchestra:0:0} & \code{aisle:orchestra:5:1}\\
\code{balcony-rim:0:0} & \code{corridor:1:1}
\end{tabular}
\end{center}
Its training clearances were 54.921, 77.165 and 79.059 seconds, respectively,
with mean 70.382 seconds. Automatic coverage had mean 72.529 seconds across
these scenarios. All training runs completed.

On held-out seed 7, the selected layout produced:
\begin{center}
\begin{tabular}{lrrrr}
\toprule
Scenario & Baseline (s) & Guided (s) & $\Delta T$ (s) & Change\\
\midrule
All open & 56.900 & 56.900 & 0.000 & 0.000\%\\
Blocked; immediate news & 64.836 & 72.565 & $+7.730$ & $+11.922\%$\\
Blocked; delayed news & 64.836 & 72.926 & $+8.090$ & $+12.478\%$\\
\bottomrule
\end{tabular}
\end{center}
The two blocked no-manager results are identical because staff-notification
timing cannot affect absent staff. All six held-out runs completed, but
guidance was slower in the blocked cases. This layout did not establish a
reliable held-out benefit under obstruction.

In the all-open held-out case, $T_{100}$ was identical while $T_{95}$ changed
from 46.957 to 47.820 seconds. This is an observed example of equal final
clearance despite a changed evacuation process. In general,
\[
 \max_i t_i^{(0)}=\max_i t_i^{(1)}
 \quad\not\Rightarrow\quad t_i^{(0)}=t_i^{(1)}\text{ for all }i.
\]

\subsection{Interpretation and limits of the results}
The implementation transmits directions, changes utility maximizers,
redistributes exits and respects physical closures. It also permits managers
to slow evacuation: shortest walking distance omits downstream congestion,
and staff occupy physical space. The observed slowdowns are compatible with
these mechanisms, but the saved summaries do not isolate a single causal
explanation for every case.

Only one training audience and one held-out audience were used in each small
search. No confidence intervals are available. These experiments do not
establish optimal placement, generally beneficial authority strength, or a
reliable percentage improvement. More audience seeds, wider placement search
and separate parameter sensitivity are needed. Congestion-aware directions
would be a different routing policy to compare with shortest-distance guidance.

\subsection{Result provenance and implementation verification}
The write-up reads saved artifacts and introduces no new simulations.
Ordinary-policy results and local decision diagnostics come from
\path{results/crowd-gui-timing-investigation.json}. Paired initial-layout
results are in \path{results/crowd-mode-verification.json}, with per-person
exit data in the corresponding \code{crowd-check-*-300.json} files.
Open search results are in \path{results/crowd-search-smoke-300.json}; the
three-scenario screening is in \path{results/crowd-scenario-smoke-300.json}.
Reusable selected-layout files are saved beside those reports. Settings are
in \path{crowd_manager_settings.json}, \path{crowd_manager_scenarios.json}
and \path{bailey_motion/default_policy_settings.json}.

Final verification records 25 focused headless tests and one real hidden-Tk
test, all passing, with 300 audience members in every real engine. An exact
pre/post baseline snapshot comparison passed. Checks also covered first-floor
exterior crossings, no passage through the scheduled closed threshold after
closure, matched initialization and finite exports. These verify implementation
behavior, not empirical human response. See
\path{results/crowd-final-verification.json}.

\section{Research provenance and limits}
The simulation--genetic-search architecture adapts von Schantz and Ehtamo,
\href{https://doi.org/10.1016/j.physa.2022.127011}{Physica A 594 (2022), 127011}.
Their moving-guide rules are not imported as stationary-manager behavior.
Clustering/distance-dependent allocation in Yang et al.,
\href{https://doi.org/10.1016/j.amc.2019.124711}{Applied Mathematics and Computation 365 (2020), 124711},
motivates the coverage initializer and distance sensitivity. The coverage
initializer itself is our chosen adaptation.

Chu et al., \href{https://doi.org/10.1016/j.trc.2017.10.009}{Transportation Research C 85 (2017), 664--683},
motivates explicit compliance and changing route information; this implementation
does not copy its fitted logistic model. CellEVAC,
\href{https://doi.org/10.1016/j.ssci.2021.105215}{Safety Science 139 (2021), 105215},
provides a related decision/movement and simulation-optimization architecture.
NIST \href{https://doi.org/10.6028/NIST.TN.1822}{TN 1822} motivates separate
checks of navigation, exit usage, route availability, flow and uncertainty.

The proximity/memory functions, stationary communication model, six-staff
constraint and parameter values are explicit project hypotheses. Experiments
on this schematic geometry cannot establish deployment suitability for an
as-built venue. More training/validation audiences, parameter sensitivity and
measured geometry are needed before interpreting a placement as robust.
\end{document}
